\begin{abstract}
Resetting or rescaling Adam's moment estimates or its bias-correction counter has been proposed as an
inexpensive remedy for the loss of learning ability under non-stationarity. Such optimizer-state
interventions also change the size of subsequent updates, so their effect is confounded with update
scale. We examine this confound in a small, fully paired synthetic setting. First, we observe the
phenomenon: on two synthetic 50-task streams and three seeds, a network that has trained on the whole
stream learns a fixed, never-seen probe task worse than the same network after its first task, with
identical probe batches, dropout masks and update budget (difference in probe learning-curve area
\RTLateMinusEarlyAuc{} and \LPLateMinusEarlyAuc{}). Second, from that late checkpoint we branch
state-only interventions with exact per-moment bias correction and compare each to a control that keeps
the original optimizer state but copies the intervention's per-step, per-tensor update norms. No
correctly bias-corrected intervention exceeded its update-norm control by our pre-declared minimum
interesting effect of $0.02$. Replacing the optimizer state by a fresh one did not recover the gap;
resetting only the counter is, with $\eps=0$, a per-tensor learning-rate schedule; and resetting only the
second moment produced very large updates despite exact bias correction. The largest gains came from a
deliberately mis-corrected reset with several-fold larger updates, for which no learning-rate-matched
control was run. These are observations under oracle intervention timing, one probe task per condition,
three previously inspected seeds and no real data; they do not show that optimizer-state repair cannot help.
\end{abstract}
